\scanpage{29}
\mspage{3}

déterminée par la condition que l'aire totale de $X$ soit $=1$. On trouve une \uline{métrique plate} caractérisée par la condition d'être invariante par translation, et d'aire \deleted{\unsure{volume}} totale égale à $1$ -- donc stable par tt autom conforme (pas seulemt holomorphe). Il s'ensuit que dans le cas non orienté, la métrique canonique de $X^{\mathrm{or}}$ (qui est une courbe elliptique $/\mathbb C$) se descend à $X^{\unsure{\mathrm{or}|n}}$.

$\deleted{\unsure{3}}2'$) Cas $X\simeq\mathbb C^*$, ou \inserted{$X\simeq$} plan projectif réel privé d'un point [i.e $X$ non orientable et $X^{\mathrm{or}}\simeq\mathbb C^*$]. Comme dans 2, dans le cas $X\simeq\mathbb C^*$, \deleted{plus intrinsèquement $X\simeq$} \deleted{\unsure{de|par}}par descente à partir du rev. universel on trouve une métrique canonique mod multiplication par un scalaire $>0$, \uline{métrique plate} également. \deleted{\unsure{Comme}} les autom de $\mathbb C^*$ \inserted{sont} \deleted{\unsure{provient|peuvent}}\inserted{\unsure{soit|être|hol}} de la forme $z\longrightarrow az$, soit $z\mathrel{\unsure{-|\to}}a/z$, qui conservent la métrique [car sur $\widetilde{\deleted{\mathbb C}X}\simeq\mathbb C$ cela correspond à des autom $z\mathrel{\unsure{-|\to}}z+\alpha$, $z\mathrel{\unsure{\vdash\!\dashv|\mapsto}}-z+\alpha$], Donc en divisant la métrique sur $\mathbb C^*$, \deleted{de façon}\inserted{\unsure{en|on}}
% archive sheet number: 31
